
\section{Introduction}
\label{sec:introduction}




Over the last decade, lattices have emerged as a very attractive
foundation for cryptography.  The appeal of lattice-based primitives
stems from the fact that their security can be based on
\emph{worst-case} hardness assumptions, that they appear to remain
secure even against \emph{quantum} computers, that they can be
quite efficient, and that, somewhat surprisingly, for certain advanced
tasks such as fully homomorphic encryption no other cryptographic
assumption is known to suffice.

Virtually all recent lattice-based cryptographic schemes are based
directly upon one of two natural average-case problems that have
been shown to enjoy worst-case hardness guarantees: the
\emph{short integer solution} ($\sis$) problem and the
\emph{learning with errors} ($\lwe$) problem.
The former dates back to Ajtai's groundbreaking
work~\cite{Ajtai96}, who showed that it is at least as
hard as approximating several worst-case lattice problems, such as the
(decision version of the) shortest vector problem, known as $\gapsvp$, to within a
polynomial factor in the lattice dimension. This hardness result
was tightened in followup work (e.g.,~\cite{DBLP:journals/siamcomp/MicciancioR07}),
leading to a somewhat satisfactory understanding of the hardness of the $\sis$ problem.
The $\sis$ problem
has been the foundation for
one-way~\cite{Ajtai96} and
collision-resistant hash
functions~\cite{DBLP:journals/eccc/ECCC-TR96-042}, identification
schemes~\cite{DBLP:conf/crypto/MicciancioV03,DBLP:conf/pkc/Lyubashevsky08,DBLP:conf/asiacrypt/KawachiTX08},
and digital
signatures~\cite{DBLP:conf/stoc/GentryPV08,DBLP:conf/eurocrypt/CashHKP10,DBLP:conf/pkc/Boyen10,DBLP:conf/eurocrypt/MicciancioP12,DBLP:conf/eurocrypt/Lyubashevsky12}. 

Our focus in this paper is on the latter problem, learning with errors.
In this problem our goal is to distinguish with some non-negligible advantage between the following
two distributions:
\[
\parens{\parens{ \veca_i, \inner{\vc{a}_i, \vc{s}}+e_i \bmod q}}_{i}
\quad \text{and} \quad
\parens{\parens{ \veca_i, u_i }}_{i}~,
\]
where $\vecs$ is chosen uniformly from $\Z_q^n$
and so are the $\veca_i \in \Z_q^n$,  $u_i$ are chosen uniformly from $\Z_q$,
and the ``noise'' $e_i \in \Z$ is sampled from some distribution supported on small numbers, typically a (discrete) Gaussian
distribution with standard deviation $\alpha q$ for $\alpha = o(1)$.
%\dnote{I replaced $\ll$ by $o(\cdot)$}


The $\lwe$ problem has proved to be amazingly versatile, serving as
the basis for a multitude of cryptographic constructions:
secure public-key encryption under both
chosen-plaintext~\cite{DBLP:journals/jacm/Regev09,DBLP:conf/crypto/PeikertVW08,DBLP:conf/ctrsa/LindnerP11}
and
chosen-ciphertext~\cite{DBLP:conf/stoc/PeikertW08,DBLP:conf/stoc/Peikert09,DBLP:conf/eurocrypt/MicciancioP12}
attacks, oblivious transfer~\cite{DBLP:conf/crypto/PeikertVW08},
identity-based
encryption~\cite{DBLP:conf/stoc/GentryPV08,DBLP:conf/eurocrypt/CashHKP10,DBLP:conf/eurocrypt/AgrawalBB10,DBLP:conf/crypto/AgrawalBB10},
various forms of leakage-resilient cryptography (e.g.,
\cite{DBLP:conf/tcc/AkaviaGV09,DBLP:conf/crypto/ApplebaumCPS09,DBLP:conf/innovations/GoldwasserKPV10}),
fully homomorphic
encryption~\cite{DBLP:conf/focs/BrakerskiV11,DBLP:conf/innovations/BrakerskiGV12,DBLP:conf/crypto/Brakerski12}
(following the seminal work of Gentry~\cite{DBLP:conf/stoc/Gentry09}),
and much more. It was also used to show hardness of learning problems~\cite{DBLP:conf/focs/KlivansS06}.

Contrary to the $\sis$ problem, however, the hardness of $\lwe$ is not
sufficiently well understood. The main hardness reduction for
$\lwe$~\cite{DBLP:journals/jacm/Regev09} is similar to the one for $\sis$
mentioned above, except that it is \emph{quantum}. This means that
the existence of an efficient algorithm for $\lwe$, even
a classical (i.e., non-quantum) one, only implies the existence
of an efficient \emph{quantum} algorithm for lattice problems.
This state of affairs is quite unsatisfactory: even though
one might conjecture that efficient quantum algorithms for lattice
problems do not exist, our understanding of quantum algorithms
is still at its infancy. It is therefore highly desirable to
come up with a \emph{classical} hardness reduction for $\lwe$.

Progress in this direction was made
by~\cite{DBLP:conf/stoc/Peikert09} (with some simplifications in the
followup by Lyubashevsky and
Micciancio~\cite{DBLP:conf/crypto/LyubashevskyM09}).  The main result
there is that $\lwe$ with \emph{exponential} modulus is as hard as
some standard lattice problems using a
classical reduction.  
%\dnote{I explained which problems in the previous sentence} 
As that hardness result crucially relies on the
exponential modulus, the open question remained as to whether $\lwe$
is hard for smaller moduli, in particular polynomial moduli. In
addition to being an interesting question in its own right, this question is of special importance since many cryptographic applications, as well as the learning theory
result of Klivans and Sherstov~\cite{DBLP:conf/focs/KlivansS06}, are instantiated in this setting.
%In
%addition to being an interesting question in its own right, we note
%that most applications use a polynomial modulus; although many
%applications can be extended to deal with an exponential modulus, a
%polynomial modulus is sometimes necessary, say in the learning theory
%result of Klivans and Sherstov~\cite{DBLP:conf/focs/KlivansS06}.
%\onote{add the ut ot?}
Some additional evidence that reducing the
modulus is a fundamental question comes from the Learning Parity with
Noise (LPN) problem, which can be seen as $\lwe$ with modulus $2$
(albeit with a different error distribution), and whose hardness is a
long-standing open question.  We remark
that~\cite{DBLP:conf/stoc/Peikert09} does include a classical hardness
of $\lwe$ with polynomial modulus, albeit one based on a
non-standard lattice problem, whose hardness is arguably as debatable
as that of the $\lwe$ problem itself.


To summarize, prior to our work, the existence of
an efficient algorithm for $\lwe$ with polynomial modulus was only known to imply an efficient \emph{quantum}
algorithm for lattice problems, or an efficient classical algorithm for a
non-standard lattice problem. While both consequences are unlikely,
they are arguably not as earth-shattering as an efficient classical
algorithm for lattice problems. Hence, some concern about the hardness of
$\lwe$ persisted, tainting the plethora of cryptographic applications
based on it.

\myparagraph{Main result}
We provide the first classical hardness reduction of $\lwe$ with polynomial modulus. Our reduction is
the first to show that the existence of an efficient classical algorithm for $\lwe$
with any subexponential modulus would indeed have earth-shattering consequences: it would imply an efficient algorithm
for worst-case instances of standard lattice problems.

\onote{should we add somewhere a formal statement of our main result? at least a combination of sections 3 and 4?}
\begin{theorem}[Informal]\label{thm:informal}
Solving $n$-dimensional $\lwe$ with $\poly(n)$ modulus
implies an equally efficient solution to a worst-case lattice problem in dimension $\sqrt{n}$.
\end{theorem}
%
As a result, we establish the hardness of all known applications of polynomial-modulus $\lwe$ based on
classical worst-case lattice problems, previously only known under a quantum assumption.

\myparagraph{Techniques}
Even though our main theorem has the flavor of a statement in computational complexity,
its proof crucially relies on a host of ideas coming from recent progress in cryptography,
most notably recent breakthroughs in the construction of fully homomorphic encryption schemes.

At a high level, our main theorem is a ``modulus reduction'' result:
we show a reduction from $\lwe$ with large modulus $q$ and dimension $n$ to
$\lwe$ with (small) modulus $p=\poly(n)$ and dimension $n \log_2 q$. Theorem~\ref{thm:informal}
now follows from the main result in~\cite{DBLP:conf/stoc/Peikert09}, which shows that
the former problem with $q=2^n$ is as hard as $n$-dimensional $\gapsvp$.
We note that the increase in dimension from $n$ to $n \log_2 q$ is to be expected,
as it essentially preserves the number of possible secrets (and hence the running time
of the naive brute-force algorithm).%invariant

Very roughly speaking,
the main idea in modulus reduction is to map $\Z_q$ into $\Z_p$ through the naive mapping
that sends any $a \in \{0,\ldots,q-1\}$ to $\floor{pa/q} \in \{0,\ldots,p-1\}$.
This basic idea is confounded by two issues. The first is that if carried out naively, this
transformation introduces rounding artifacts into $\lwe$, ruining the distribution
of the output. We resolve this issue by using a more careful Gaussian randomized rounding procedure (Section~\ref{sec:modexpansion}).
A second serious issue is that in order for the rounding errors not to be
amplified when multiplied by the $\lwe$ secret $\vecs$, it is essential to assume that $\vecs$
has small coordinates. A major part of our reduction (Section~\ref{sec:shortsecret}) is therefore dedicated to showing
a reduction from $\lwe$ (in dimension $n$) with arbitrary secret in $\Z_q^n$ to $\lwe$ (in dimension $n \log_2 q$)
with a secret chosen uniformly over $\{0,1\}$. This follows from a careful hybrid argument (Section~\ref{sec:theshortsecretreduction}) combined with a
hardness reduction to the so-called ``extended-$\lwe$'' problem, which is a variant of $\lwe$ in which we have some control
over the error vector (Section~\ref{sec:extlwe}).


We stress that even though our proof is inspired by and has analogues
in the cryptographic literature, the details of the reductions are very different.
In particular, the idea of modulus reduction plays a key role in recent work on fully homomorphic
encryption schemes, giving a way to control the noise growth during homomorphic
operations~\cite{DBLP:conf/focs/BrakerskiV11,DBLP:conf/innovations/BrakerskiGV12,DBLP:conf/crypto/Brakerski12}.
However, since the goal there is merely to preserve the functionality of the scheme,
their modulus reduction can be performed in a rather naive way similar to the one outlined above,
and so the output of their procedure does not constitute a valid $\lwe$ instance.
In our reduction we need to perform a much more delicate modulus reduction,
which we do using Gaussian randomized rounding, as mentioned above.


%
The idea of reducing $\lwe$ to have a $\{0,1\}$ secret also exists already in the cryptographic literature:
precisely such a reduction was shown by 
Goldwasser et al.~\cite{DBLP:conf/innovations/GoldwasserKPV10}
who were motivated by questions in leakage-resilient cryptography.
Their reduction, however, incurred a severe blow-up in the noise rate, making it useless for our purposes. 
In more detail, not being able to
faithfully reproduce the $\lwe$ distribution in the output, they resort to 
hiding the faults in the output distribution under a huge independent 
fresh noise, in order to make it close to the correct one. The trouble 
with this ``noise flooding'' approach
%\dnote{I reformulated these sentences to make ``noise flooding'' less 
% mysterious}
%what has been called ``noise flooding'' in order to hide the faults
%in the output distribution, and to make it close to the correct one. \znote{Maybe say what noise flooding is? It seems somewhat mysterious like this.}
%The trouble with this approach 
is that the amount of noise
one has to add depends on the running time of the algorithm solving
the target $\{0,1\}$-$\lwe$ problem, which in turn forces the modulus
to be equally big.
So while in principle we could use the reduction
from~\cite{DBLP:conf/innovations/GoldwasserKPV10} (and shorten our proof
by about a half), this would lead to a qualitatively much weaker result:
the modulus and the approximation ratio for the worst-case lattice problem
would both grow with the running time of the $\{0,1\}$-$\lwe$ algorithm.
In particular, we would not be able to show that for some fixed
polynomial modulus, $\lwe$ is a hard problem; instead, in order to capture all polynomial time algorithms,
we would have to take a super-polynomial
modulus, and rely on the hardness of worst-case lattice problem to within super-polynomial
approximation factors.
In contrast, with our reduction,
the modulus and the approximation ratio both remain fixed independently of the target $\{0,1\}$-$\lwe$ algorithm.


As mentioned above, our alternative to the reduction
in~\cite{DBLP:conf/innovations/GoldwasserKPV10} is based
on a hybrid argument combined with a new hardness reduction
for the ``extended LWE'' problem, which is a variant of $\lwe$ in which in addition to the LWE samples, we also get to see the inner product
of the vector of error terms with a vector $\vecz$ of our 
%\dnote{Shall we add (``adversorial'')?}\onote{i prefer to keep as is} 
choosing. This problem has
its origins in the cryptographic literature, namely
in the work of O'Neill, Peikert, and Waters~\cite{DBLP:conf/crypto/ONeillPW11} on (bi)deniable encryption
and the later work of Alperin-Sheriff and Peikert~\cite{DBLP:conf/pkc/Alperin-SheriffP12}
on key-dependent message security. The hardness reductions included in those
papers are not sufficient for our purposes, as they cannot handle
large moduli or error terms, which is crucial in our setting. We therefore
provide an alternative reduction which is conceptually much simpler, and
essentially subsumes both previous reductions.
Our reduction works equally well with
exponential moduli and correspondingly long error vectors, a case earlier reductions could not handle.



\myparagraph{Broader perspective}
As a byproduct of the proof of Theorem~\ref{thm:informal}, we obtain several results that
shed new light on the hardness of $\lwe$. Most notably, our modulus reduction result
in Section~\ref{sec:modexpansion} is actually far more general, and can be used to show
a ``modulus expansion/dimension reduction'' tradeoff. Namely, it shows a reduction
from $\lwe$ in dimension $n$ and modulus $p$ to $\lwe$ in dimension $n/k$ and modulus $p^k$ (see Corollary~\ref{cor:mod-dim-tradeoff}).
Combined with our modulus reduction, this has the following interesting consequence:
the hardness of $n$-dimensional $\lwe$ with modulus $q$ is a function of the quantity $n \log_2 q$.
In other words, varying $n$ and $q$ individually while keeping $n \log_2 q$ fixed
essentially preserves the hardness of $\lwe$.

Although we find this statement quite
natural (since $n \log_2 q$ represents the number of bits in the secret), it has
some surprising consequences. One is that $n$-dimensional $\lwe$ with modulus $2^n$
is essentially as hard as $n^2$-dimensional $\lwe$ with polynomial modulus.
As a result, $n$-dimensional $\lwe$ with modulus $2^n$, which was shown
in~\cite{DBLP:conf/stoc/Peikert09} to be as hard as $n$-dimensional lattice
problems using a classical reduction, is actually as hard as $n^2$-dimensional
lattice problems using a quantum reduction. The latter is presumably a much
harder problem, requiring $\exp(\widetilde\Omega(n^2))$ time to solve.
This corollary highlights an inherent quadratic loss in the classical reduction
of~\cite{DBLP:conf/stoc/Peikert09} (and as a result also our Theorem~\ref{thm:informal})
compared to the quantum one in~\cite{DBLP:journals/jacm/Regev09}.

\onote{for the final version, consider adding precise statements for this and the next point}
A second interesting consequence is that $1$-dimensional $\lwe$ with
modulus $2^n$ is essentially as hard as $n$-dimensional $\lwe$ with
polynomial modulus. The $1$-dimensional version of $\lwe$ is closely
related to the Hidden Number Problem of Boneh and
Venkatesan~\cite{DBLP:conf/crypto/BonehV96}. It is also essentially equivalent to the
Ajtai-Dwork-type~\cite{DBLP:conf/stoc/AjtaiD97} cryptosystem in~\cite{DBLP:journals/jacm/Regev04}, as follows from simple reductions
similar to the one in the appendix of~\cite{RegevSurvey}.
Moreover, the $1$-dimensional version can be seen as a special case of
the Ring-$\lwe$ problem introduced
in~\cite{DBLP:conf/eurocrypt/LyubashevskyPR10} (for ring dimension~1,
i.e., ring equal to~$\Z$). This allows us, via the ring switching
technique from~\cite{DBLP:conf/scn/GentryHPS12}, to obtain the first
hardness proof of Ring-$\lwe$, with arbitrary ring dimension and
exponential modulus, under the hardness of problems on general lattices (as
opposed to just ideal lattice problems). In addition, this leads to the
first hardness proof for the Ring-$\sis$
problem~\cite{DBLP:conf/icalp/LyubashevskyM06,DBLP:conf/tcc/PeikertR06}
with exponential modulus under the hardness of general lattice problems, via the standard
$\lwe$-to-$\sis$ reduction. (We note that since both results are obtained by
scaling up from a ring of dimension $1$, the hardness does not improve
as the ring dimension increases.)

A final interesting consequence of our reductions is that (the
decision form of) \lwe is hard with an arbitrary huge modulus, e.g., a
prime; see Corollary~\ref{cor:mod-reduction-2}.  Previous results
(e.g.,~\cite{DBLP:journals/jacm/Regev09,DBLP:conf/stoc/Peikert09,DBLP:conf/crypto/MicciancioM11,DBLP:conf/eurocrypt/MicciancioP12})
required the modulus to be \emph{smooth}, i.e., all its prime divisors
had to be polynomially bounded.

\onote{for final version: decide what to do with real-LWE}

\onote{for final version, consider applying these ideas to SIS, especially the 1-dim case which is like subset sum, and more generally, modulus/dimension tradeoff for SIS. Chris said:
The idea is to do a similar transformation as in Lemma 4.4, but on SIS samples, going from a in $\Z_q^n$ to $a' \in \Z_{q'}^{n'}$.  Then given a good SIS solution $x$ to $A'x = 0$, we also have $G^T A' x = 0$ and therefore $A x \approx 0$.  Thus $x$ is a solution to the original SIS instance, because it suffices to find a preimage of "nearly" 0.  This is because if $Ax = e \approx 0$, then we get the homogeneous solution $[I | A] (e; x) = 0$, and the instance $[I | A]$ is just SIS in (hermite) normal form.}

\dnote{Something that I find interesting is that the complexity of the
  best known attack is~$2^{O(n \log q/ \log^2 \alpha)}$, which is
  constant wrt $n \log q$. This was not reflected by our previous
  understanding of the complexity of \lwe, and we manage to close the
  gap between the hardness results and the algorithms. By the way,
  this~$2^{O(n \log q/ \log^2 \alpha)}$ bound suggests it may be
  possible to transfer some of~$n$ and~$\log q$ into~$\log \alpha$. Do
  we know \lwe self-reductions from, say, $q,\alpha$ to~$p,\beta$ such
  that~$\log q/ |\log \alpha| = \log p/|\log \beta|$? or from
  $n,\alpha$ to~$n',\beta$ such that~$n/ |\log \alpha| = n'/|\log
  \beta|$}

\onote{maybe mention Ajtai's 2005 paper using Dirichlet's problem?}

\myparagraph{Open questions}
As mentioned above, our Theorem~\ref{thm:informal} inherits from~\cite{DBLP:conf/stoc/Peikert09}
a quadratic loss in the dimension, which does not exist in the quantum reduction~\cite{DBLP:journals/jacm/Regev09}
nor in the known hardness reductions for $\sis$. 
At a technical level, this quadratic loss stems from the fact that the reduction in~\cite{DBLP:conf/stoc/Peikert09} is not iterative. In contrast,
the quantum reduction in~\cite{DBLP:journals/jacm/Regev09} as well as the reductions for $\sis$ are iterative, and as a result do not incur
the quadratic loss. We note that an additional side effect of the non-iterative reduction is that the hardness in Theorem~\ref{thm:informal} 
and~\cite{DBLP:conf/stoc/Peikert09} is based only on the worst-case lattice problem \gapsvp (and the essentially equivalent \bdd and \usvp~\cite{DBLP:conf/crypto/LyubashevskyM09}), and not on problems like \sivp, which the quantum reduction of~\cite{DBLP:journals/jacm/Regev09} and the hardness
reductions for \sis can handle. One case where this is very significant is when dealing with ideal lattices, as in the hardness 
reduction for Ring-\lwe, since \gapsvp turns out to be an easy problem there. 

We therefore believe that it is important to understand whether there exists a classical
reduction that does not incur the quadratic loss inherent in~\cite{DBLP:conf/stoc/Peikert09}
and in Theorem~\ref{thm:informal}.
In other words, is $n$-dimensional $\lwe$ with polynomial modulus classically as hard as $n$-dimensional
lattice problems (as opposed to $\sqrt{n}$-dimensional)? This would constitute
the first full dequantization of the quantum reduction in~\cite{DBLP:journals/jacm/Regev09}.

While it is natural to conjecture that the answer to this question is positive,
a negative answer would be quite tantalizing. In particular, it is
conceivable that there exists a (classical) algorithm for $\lwe$ with polynomial modulus running in time
$2^{O(\sqrt{n})}$. Due to the quadratic expansion in Theorem~\ref{thm:informal}, this
would not lead to a faster classical algorithm for lattice problems; it would,
however, lead to a $2^{O(\sqrt{n})}$-time \emph{quantum} algorithm for
lattice problems using the reduction in~\cite{DBLP:journals/jacm/Regev09}.
The latter would be a major progress in quantum algorithms, yet is not
entirely unreasonable; in fact, a $2^{O(\sqrt{n})}$-time quantum algorithm
for a somewhat related quantum task was discovered by
Kuperberg~\cite{DBLP:journals/siamcomp/Kuperberg05} (see also~\cite{DBLP:journals/siamcomp/Regev04}).




%%% Local Variables:
%%% mode: latex
%%% TeX-master: "lwehard"
%%% End:
